% Part III-C, sections that do not depend on the results (included by main_en.tex).

\section{Judging unfinished solutions}\label{sec:prefix}

Tree rollouts~\cite{hou2025treerl} sample several prefixes of a solution, continue some of them and share each prefix
among the continuations. A judge that can tell early which prefixes lead to correct answers could decide which ones to
continue. This saves the cost of continuing dead ends and concentrates samples where the policy can succeed. We first
measure offline whether a 4B judge can make that call, then train with it.

\paragraph{Prefix study.} We use GRPO's validation problems (500 MATH training problems held out from training,
\S2 of Part II). The base policy samples 4 solutions per problem as in GRPO (temperature 1, the training prompt). Every
solution is cut at 128, 256 and 512 tokens if it runs that long. Each prefix gets 4 continuations to the training
length (1024 tokens in all), and the verifier checks them. The fraction of a prefix's continuations that end correct is
its Monte-Carlo value. The signals are:
\begin{itemize}
\item \kev{}'s $p(\text{yes})$ for \emph{on track} (``Is this unfinished solution on track to reach the correct final
  answer\ldots?'');
\item \kev{}'s $p(\text{yes})$ for Part II's \emph{correct} question, with the prefix as the proposed solution;
\item the AR judge's $p(\text{yes})$ for \emph{on track};
\item the policy's own mean token log-probability of the prefix, which costs nothing.
\end{itemize}
For each cut length we report each signal's AUROC for a continuation ending correct. It is reported pooled and within
problems, the comparison that matters when choosing among one problem's prefixes. We also report each signal's
Spearman correlation with the Monte-Carlo value.

\paragraph{Pruning experiment.} On the problems whose 4 solutions all run past the cut, a signal keeps the 2 prefixes
it ranks highest. We compare the mean value of the kept prefixes with the mean of all 4, which is what a random choice
keeps in expectation. We also compare it with an oracle that ranks by half of the continuations and is scored on the
other half. There are 95\% bootstrap intervals over problems. We also evaluate the rule of the withdrawn v3 paper:
prune a prefix when the judge's $p(\text{yes})<0.3$. For it we report how much it prunes, how many pruned prefixes were
dead (value 0), and how many correct continuations it throws away.

\section{Tree GRPO}\label{sec:tree}

Each prompt's group of 8 leaves grows from $R=4$ sampled roots of 128 tokens. $K=2$ of them are continued, with 4
samples each, from the exact token prefix. \texttt{tree\_kev} keeps the two roots with the highest \kev{}
\emph{on-track} score; \texttt{tree\_random} keeps two at random. In distribution, \texttt{tree\_random} equals sampling
two roots only. A root that has already ended is a complete solution, and if kept its 4 leaves are copies of it. Leaves
are trained like Part II's flat completions: the full root and continuation tokens, the verifier's reward, and the
group baseline over the prompt's 8 leaves. Everything else is Part II's \texttt{rlvr} setup. Every tenth step, the roots
that were not kept are continued too, for the log only. This measures what the selection gained on that step against
keeping all roots, and against the best two in hindsight.

Two effects pull in opposite directions. Selecting roots by a judge raises the share of correct leaves when the judge
is right, which gives GRPO more groups with a learning signal. But it also makes the kept roots no longer samples of the
policy, so the gradient is no longer an unbiased estimate of the policy gradient. Comparing \texttt{tree\_kev} with
\texttt{tree\_random} (same tree shape, unbiased selection) isolates the selection, and comparing \texttt{tree\_random}
with Part II's \texttt{rlvr} isolates the tree shape.
